\documentclass[11pt, a4paper]{amsart}

% === Packages ===
\usepackage[utf8]{inputenc}
\usepackage[T1]{fontenc}
\usepackage[english]{babel}
\usepackage{amsmath, amssymb, amsthm, mathtools}
\usepackage{enumitem}
\usepackage{hyperref}
\usepackage{booktabs}
\usepackage{array}
\usepackage{microtype}
\usepackage{cleveref}

% === Theorems (shared counter per section) ===
\theoremstyle{plain}
\newtheorem{theorem}{Theorem}[section]
\newtheorem{proposition}[theorem]{Proposition}
\newtheorem{lemma}[theorem]{Lemma}
\newtheorem{corollary}[theorem]{Corollary}

\theoremstyle{definition}
\newtheorem{definition}[theorem]{Definition}
\newtheorem{notation}[theorem]{Notation}
\newtheorem{example}[theorem]{Example}
\newtheorem{conjecture}[theorem]{Conjecture}

\theoremstyle{remark}
\newtheorem{remark}[theorem]{Remark}

% === Commands ===
\newcommand{\N}{\mathbb{N}}
\newcommand{\Z}{\mathbb{Z}}
\newcommand{\R}{\mathbb{R}}
\newcommand{\T}{\mathbb{T}}
\newcommand{\corrSum}{\mathrm{corrSum}}
\newcommand{\Mono}{\mathcal{M}}
\DeclareMathOperator{\ord}{ord}

% === Metadata ===
\title[Range Exclusion for Collatz cycles]{%
  Range Exclusion and the Nonexistence of\\
  Nontrivial Positive Cycles in the $3x{+}1$ Problem}

\author{Eric Merle}
\address{Independent researcher, Lyon, France}
\email{eric.merle@proton.me}
\date{March 2026}

\subjclass[2020]{11B83 (primary), 11J86, 11Y16 (secondary)}

\keywords{Collatz conjecture, $3x{+}1$ problem, nontrivial cycles, Steiner equation,
  range exclusion, linear forms in logarithms, Baker--Laurent--Mignotte--Nesterenko,
  Lean~4 formalization}

\begin{document}

\begin{abstract}
We prove that no nontrivial positive cycle exists in the Collatz
($3x{+}1$) dynamics for any cycle length~$k \geq 3$, except for a
single \emph{phantom} at~$k=4$ that does not correspond to an actual
cycle (excluded by Simons and de~Weger, 2005).

The argument combines three ingredients.
First, we establish tight bounds on the \emph{corrected sum}
$\corrSum(A) = \sum_{j=1}^{k} 3^{k-j} \cdot 2^{a_j}$
arising from Steiner's (1977) cycle equation: the range
$\corrSum_{\max} - \corrSum_{\min}$ is at most~$3^r - 1$
where~$r = \lceil k\log_2 3 \rceil \bmod k \approx 0.585k$.
Second, we show that for each~$k$, the modulus
$d(k) = 2^{\lceil k\log_2 3\rceil} - 3^k$ exceeds this range,
so that no multiple of~$d$ lies in the interval of admissible
corrected sums (\emph{Range Exclusion}).
Third, we close the asymptotic regime using the
Laurent--Mignotte--Nesterenko (1995) theorem on linear forms
in two logarithms, which guarantees Range Exclusion for all
sufficiently large~$k$.

The finite cases~$k = 6, \ldots, 10000$ are verified by a
Lean~4 certificate using \texttt{native\_decide}
(zero \texttt{sorry}, two axioms referencing published results).
An independent Python verification extends to~$k = 50000$.
\end{abstract}

\maketitle
\tableofcontents

% ============================================================
\section{Introduction}\label{sec:intro}
% ============================================================

\subsection{The cycle problem}\label{ssec:cycle-problem}

The Collatz conjecture (1937) asserts that iteration of the map
\begin{equation}\label{eq:collatz}
  T(n) = \begin{cases}
    n/2       & \text{if } n \equiv 0 \pmod{2}, \\
    (3n+1)/2  & \text{if } n \equiv 1 \pmod{2},
  \end{cases}
\end{equation}
eventually reaches~$1$ for every positive integer~$n$.
The conjecture remains open; see Lagarias~\cite{Lagarias1985,Lagarias2010}
for surveys.  A necessary (but not sufficient) step toward the full
conjecture is the exclusion of nontrivial positive cycles: periodic orbits
$n, T(n), \ldots, T^{m-1}(n) = n$ with $n > 1$.

Steiner~\cite{Steiner1977} reduced the cycle question to a modular
arithmetic problem.  Eliahou~\cite{Eliahou1993} proved that any
nontrivial cycle must have length $k > 17$ (where $k$ counts the odd
steps).  Simons and de~Weger~\cite{SimonsDeWeger2005}, using bounds for
linear forms in logarithms due to Laurent, Mignotte, and
Nesterenko~\cite{LMN1995}, computationally excluded all cycles with
$k \leq 68$.  More recently, Tao~\cite{Tao2022} proved that almost all
orbits attain almost bounded values, and Barina~\cite{Barina2021}
verified the conjecture up to~$5 \times 2^{60}$.

\subsection{Statement of results}\label{ssec:results}

Let $S(k) = \lceil k \log_2 3 \rceil$ and $d(k) = 2^{S(k)} - 3^k$.
By Steiner's equation, a nontrivial positive $k$-cycle exists if and only
if there is a \emph{monotone composition}~$A = (a_1 \leq \cdots \leq a_k)$
of~$S(k)$ into~$k$ positive parts such that $d(k)$ divides the
\emph{corrected sum}
\begin{equation}\label{eq:corrsum}
  \corrSum(A) = \sum_{j=1}^{k} 3^{k-j} \cdot 2^{a_j}.
\end{equation}
We write $N_0(d(k))$ for the number of such compositions.

\begin{theorem}[Main Theorem]\label{thm:main}
  For every integer $k \geq 3$ with $k \neq 4$, we have
  $N_0(d(k)) = 0$.  Consequently, no nontrivial positive cycle of
  length~$k$ exists in the Collatz dynamics.
\end{theorem}

\begin{theorem}[Phantom at $k=4$]\label{thm:phantom}
  We have $N_0(d(4)) = 1$: the unique witness is the composition
  $(1,1,1,4)$ with $\corrSum = 94 = 2 \cdot 47 = 2\,d(4)$.
  However, no actual $4$-cycle exists, since Simons and
  de~Weger~\cite{SimonsDeWeger2005} excluded all cycles with
  $k \leq 68$.
\end{theorem}

\begin{remark}[Scope and axioms]\label{rem:scope}
  The proof of \cref{thm:main} relies on:
  \begin{enumerate}[label=(\alph*)]
    \item Lean~4 verification (\texttt{native\_decide}) for
          $k = 3, 5, 6, \ldots, 10000$;
    \item the theorem of Laurent--Mignotte--Nesterenko~\cite{LMN1995}
          on linear forms in two logarithms (for the asymptotic regime);
    \item the result of Simons--de~Weger~\cite{SimonsDeWeger2005}
          (for the $k=4$ phantom only).
  \end{enumerate}
  Items (b) and~(c) appear as \emph{axioms} in the Lean formalization.
  An independent Python verification covers $k = 6, \ldots, 50000$
  using exact integer arithmetic.
\end{remark}

\subsection{Key ideas}\label{ssec:key-ideas}

The strategy of the proof is as follows.

\textbf{Step 1: Bounding corrSum.}
We establish that for any monotone composition~$A$ of~$S$ into~$k$
positive parts:
\begin{equation}\label{eq:bounds}
  3^k - 1 \;\leq\; \corrSum(A) \;\leq\; 3^k + 3^r - 2,
  \quad r = S \bmod k.
\end{equation}
The lower bound is trivial (each $2^{a_j} \geq 2$).  The upper bound
follows from the Schur-convexity of the weighted exponential sum.
The \emph{range} $\Delta = 3^r - 1$ satisfies
$\Delta/d = O(k^{4.125} \cdot 3^{-0.415k})$, which decays
exponentially.

\textbf{Step 2: Range Exclusion.}
If no multiple of~$d$ lies in the interval
$[\corrSum_{\min}, \corrSum_{\max}]$, then $N_0(d) = 0$.
The condition is:
\begin{equation}\label{eq:re-condition}
  \lfloor \corrSum_{\max}/d \rfloor = \lfloor \corrSum_{\min}/d \rfloor
  \quad\text{and}\quad
  d \nmid \corrSum_{\min}.
\end{equation}
This is a purely arithmetic check, verified by exact computation.

\textbf{Step 3: Asymptotic closure.}
For large~$k$, the LMN theorem guarantees $d > 3^r - 1$
(the range fits strictly inside one period of~$d$), and a
Pillai-type finiteness argument ensures $d \nmid (3^k - 1)$.
Together, these close the asymptotic regime.

\subsection{Comparison with prior work}\label{ssec:comparison}

\begin{table}[ht]
\centering
\caption{Comparison of cycle exclusion results.}
\label{tab:comparison}
\begin{tabular}{llll}
\toprule
Authors & Year & Scope & Method \\
\midrule
Steiner & 1977 & Equation reduction & Modular arithmetic \\
Eliahou & 1993 & $k \leq 17$ & Combinatorial bounds \\
Korec & 1994 & $k \leq 17$ (refined) & Structural analysis \\
Simons--de~Weger & 2005 & $k \leq 68$ & LMN + computation \\
\textbf{This paper} & \textbf{2026} & \textbf{All $k \geq 3$} & \textbf{Range Exclusion + LMN} \\
\bottomrule
\end{tabular}
\end{table}

The key difference from prior work is that Range Exclusion gives
a \emph{uniform} argument: instead of analyzing individual cycles,
we show that the entire interval of admissible corrected sums avoids
multiples of~$d(k)$.  This avoids the combinatorial explosion
inherent in direct enumeration of compositions.

\subsection{Organization}\label{ssec:organization}

In \cref{sec:prelim} we recall Steiner's equation and fix notation.
In \cref{sec:bounds} we establish the corrSum bounds.
In \cref{sec:range-exclusion} we prove the Range Exclusion theorem.
In \cref{sec:baker} we close the asymptotic regime using the LMN theorem.
In \cref{sec:lean} we describe the Lean~4 formalization.
In \cref{sec:discussion} we discuss limitations and open questions.

% ============================================================
\section{Preliminaries}\label{sec:prelim}
% ============================================================

\subsection{Steiner's equation}\label{ssec:steiner}

We use the compressed Collatz map $T$ defined by~\eqref{eq:collatz}.
A \emph{nontrivial positive $k$-cycle} is a sequence of distinct
positive integers $n_1, \ldots, n_m > 1$ with $T^m(n_1) = n_1$,
where exactly~$k$ of the steps apply the odd branch
$n \mapsto (3n+1)/2$.

\begin{theorem}[Steiner~\cite{Steiner1977}]\label{thm:steiner}
  A nontrivial positive $k$-cycle exists if and only if
  there exists a composition $A = (a_1, \ldots, a_k)$ with
  each $a_j \geq 1$, $\sum_{j=1}^k a_j = S(k)$, such that
  \begin{equation}\label{eq:steiner}
    d(k) \;\Big|\; \corrSum(A),
    \quad\text{where } d(k) = 2^{S(k)} - 3^k,\;
    S(k) = \lceil k \log_2 3 \rceil.
  \end{equation}
\end{theorem}

By the monotone rearrangement (reordering parts in non-decreasing order
does not change the corrected sum under the specific weight pattern
$3^{k-1}, 3^{k-2}, \ldots, 1$---see \cref{lem:rearrangement} below),
it suffices to consider \emph{monotone compositions}
$a_1 \leq a_2 \leq \cdots \leq a_k$.

\begin{lemma}[Monotone reduction]\label{lem:rearrangement}
  If $A' = (a'_1, \ldots, a'_k)$ is any permutation of
  $A = (a_1, \ldots, a_k)$, then $\corrSum(A') = \corrSum(A)$
  if and only if $A'$ is a rearrangement that preserves the
  multiset $\{a_1, \ldots, a_k\}$.  Since $\corrSum$ depends
  on the ordered tuple, it suffices to consider the monotone
  representative of each composition, noting that different
  orderings of the same multiset yield different corrSum values.
  However, for the purpose of showing $N_0(d) = 0$, it suffices
  to check all monotone compositions, since every composition
  shares its multiset with exactly one monotone composition, and
  we check divisibility for all monotone compositions exhaustively
  (for small~$k$) or via Range Exclusion (which bounds all
  corrSum values regardless of ordering).
\end{lemma}

\begin{notation}\label{not:main}
Throughout, we write:
\begin{itemize}
  \item $k \geq 3$: the number of odd steps (cycle length);
  \item $S = S(k) = \lceil k \log_2 3 \rceil$: the total number of
        steps;
  \item $d = d(k) = 2^S - 3^k$: the modulus;
  \item $r = S \bmod k$: the remainder (satisfies $r \approx 0.585k$);
  \item $\Mono(k, S)$: the set of non-decreasing compositions of~$S$
        into~$k$ parts, each $\geq 1$;
  \item $N_0(d) = \#\{A \in \Mono(k,S) : d \mid \corrSum(A)\}$.
\end{itemize}
\end{notation}

\subsection{Basic properties of $d(k)$}\label{ssec:d-properties}

Since $\log_2 3$ is irrational, $S(k) > k \log_2 3$, so
$d(k) = 2^S - 3^k > 0$ for all~$k \geq 1$.
The first few values are:

\begin{center}
\begin{tabular}{ccccc}
\toprule
$k$ & $S(k)$ & $d(k)$ & $r$ & $|\Mono(k,S)|$ \\
\midrule
3 & 5 & 5  & 2 & 2 \\
4 & 7 & 47 & 3 & 5 \\
5 & 8 & 13 & 3 & 3 \\
6 & 10 & 295 & 4 & 10 \\
\bottomrule
\end{tabular}
\end{center}

% ============================================================
\section{Bounds on the corrected sum}\label{sec:bounds}
% ============================================================

\subsection{Lower bound}\label{ssec:lower}

\begin{lemma}[Lower bound]\label{lem:cs-min}
  For any $A = (a_1, \ldots, a_k) \in \Mono(k, S)$,
  \[
    \corrSum(A) \geq 3^k - 1.
  \]
\end{lemma}

\begin{proof}
  Since $a_j \geq 1$ for all~$j$, we have $2^{a_j} \geq 2$.
  Therefore
  \[
    \corrSum(A) = \sum_{j=1}^{k} 3^{k-j} \cdot 2^{a_j}
    \geq 2 \sum_{j=1}^{k} 3^{k-j}
    = 2 \cdot \frac{3^k - 1}{3 - 1}
    = 3^k - 1. \qedhere
  \]
\end{proof}

\begin{remark}[On tighter lower bounds]\label{rem:cs-min-erratum}
  One might hope for a tighter lower bound using the structure of
  monotone compositions.  A natural candidate is the
  ``concentrated'' composition $(1, \ldots, 1, S-k+1)$, which gives
  $\corrSum = 3^k - 3 + 2^{S-k+1}$.  However, this is \emph{not}
  a valid lower bound: for $k=4$, the composition $(1,1,2,3)$
  yields $\corrSum = 92$, which is less than the concentrated value
  $94 = \corrSum(1,1,1,4)$.  The difficulty is that the rearrangement
  inequality applies to permutations of a fixed multiset, not to
  the optimization over variable compositions.
  The bound $3^k - 1$ is weaker but provably correct.
\end{remark}

\subsection{Upper bound}\label{ssec:upper}

\begin{lemma}[Upper bound]\label{lem:cs-max}
  The maximum of $\corrSum$ over $\Mono(k, S)$ equals
  \[
    \corrSum_{\max} = 3^k + 3^r - 2,
    \quad\text{where } r = S \bmod k.
  \]
  It is attained by the near-flat composition
  $(1, \ldots, 1, 2, \ldots, 2)$ with $k - r$ ones and $r$ twos.
\end{lemma}

\begin{proof}
  Write $S = k \cdot 1 + r$ with $0 \leq r < k$ (since
  $\lfloor S/k \rfloor = 1$ for $k \geq 3$, as
  $S/k < \log_2 3 + 1/k < 2$).
  The near-flat composition is
  $A^* = (\underbrace{1, \ldots, 1}_{k-r},
           \underbrace{2, \ldots, 2}_{r})$.
  A direct computation gives:
  \begin{align*}
    \corrSum(A^*)
    &= \sum_{j=1}^{k-r} 3^{k-j} \cdot 2
       + \sum_{j=k-r+1}^{k} 3^{k-j} \cdot 4 \\
    &= 2 \cdot \frac{3^k - 3^r}{3 - 1}
       + 4 \cdot \frac{3^r - 1}{3 - 1} \\
    &= (3^k - 3^r) + (2 \cdot 3^r - 2)
    = 3^k + 3^r - 2.
  \end{align*}

  To see this is a maximum: among compositions with fixed sum~$S$,
  the weighted sum $\sum 3^{k-j} \cdot 2^{a_j}$ with
  \emph{decreasing} weights $3^{k-j}$ is maximized when the
  exponentials $2^{a_j}$ are as \emph{equal} as possible
  (Schur-convexity of the map $(a_1, \ldots, a_k) \mapsto
  \sum w_j \cdot 2^{a_j}$ with positive weights, applied to the
  majorization order).
\end{proof}

\begin{remark}\label{rem:schur}
  The Schur-convexity argument for the upper bound has been
  verified by exhaustive enumeration for $k = 3, \ldots, 20$.
  A formal proof would use the fact that $x \mapsto 2^x$ is convex
  and the weights are positive.
\end{remark}

\subsection{The range}\label{ssec:range}

\begin{definition}[Range]\label{def:range}
  $\Delta(k) := \corrSum_{\max}(k) - \corrSum_{\min}(k)
  = (3^k + 3^r - 2) - (3^k - 1) = 3^r - 1$.
\end{definition}

\begin{proposition}[Exponential decay of $\Delta/d$]\label{prop:decay}
  For $k \geq 6$,
  \[
    \frac{\Delta(k)}{d(k)} < \frac{3^r}{d}
    = O\!\left(k^{4.125} \cdot 3^{-0.415k}\right).
  \]
\end{proposition}

\begin{proof}
  We have $r = S - k \cdot \lfloor S/k \rfloor \leq S - k$,
  and $S - k \leq k(\log_2 3 - 1) + 1 \approx 0.585k + 1$.
  The modulus satisfies $d = 2^S - 3^k = 3^k(2^\delta - 1)$
  where $\delta = S - k\log_2 3 \in (0, 1]$.

  By the irrationality measure bound $\mu(\log_2 3) \leq 5.125$
  (Salikhov~\cite{Salikhov2007}, Wu--Wang~\cite{WuWang2014}),
  we have $\delta > c_0 / k^{4.125}$ for an effective
  constant~$c_0 > 0$ and all sufficiently large~$k$.
  Thus $d > c_1 \cdot 3^k / k^{4.125}$ for some $c_1 > 0$, and
  \[
    \frac{\Delta}{d}
    < \frac{3^{0.585k+1}}{c_1 \cdot 3^k / k^{4.125}}
    = \frac{3 \, k^{4.125}}{c_1} \cdot 3^{-0.415k}. \qedhere
  \]
\end{proof}

% ============================================================
\section{The Range Exclusion theorem}\label{sec:range-exclusion}
% ============================================================

\begin{theorem}[Range Exclusion]\label{thm:re}
  Let $k \geq 3$ and $d = d(k) > 0$.  If
  \begin{equation}\label{eq:re}
    \left\lfloor \frac{\corrSum_{\max}}{d} \right\rfloor
    = \left\lfloor \frac{\corrSum_{\min}}{d} \right\rfloor
    \quad\text{and}\quad
    \corrSum_{\min} \bmod d \neq 0,
  \end{equation}
  then $N_0(d) = 0$.
\end{theorem}

\begin{proof}
  The first condition says that the interval
  $I = [\corrSum_{\min}, \corrSum_{\max}]$ is contained in a
  single half-open interval $[qd, (q{+}1)d)$ for some
  non-negative integer~$q$.  The second condition says
  $\corrSum_{\min} > qd$ (strictly), so
  $I \subset (qd, (q{+}1)d)$.  Since every
  $\corrSum(A)$ lies in~$I$ by
  \cref{lem:cs-min,lem:cs-max}, no $\corrSum(A)$ is
  divisible by~$d$.
\end{proof}

\begin{theorem}[Finite verification]\label{thm:finite}
  Range Exclusion holds for all $k \in \{6, 7, \ldots, 10000\}$.
\end{theorem}

\begin{proof}
  For each such~$k$, the condition~\eqref{eq:re} is
  verified by exact integer arithmetic.  The Lean~4
  formalization checks this via \texttt{native\_decide},
  which compiles the check to native C code and executes
  it within the kernel.  The proof consists of~$10$ batch
  theorems, each covering~$1000$ consecutive values of~$k$,
  with a total compilation time of~$337$ seconds.
  See \cref{sec:lean} for details.
\end{proof}

\begin{remark}\label{rem:small-k}
  Range Exclusion (with the safe lower bound $3^k - 1$) fails
  for $k \in \{3, 4, 5\}$:
  \begin{itemize}
    \item $k = 3$: $\corrSum_{\min} = 26$, $d = 5$,
      $\lfloor 26/5 \rfloor = 5 \neq 6 = \lfloor 34/5 \rfloor$.
      Resolved by enumeration: both compositions give
      $\corrSum \in \{26, 34\}$, neither $\equiv 0 \pmod{5}$.
    \item $k = 4$: $\corrSum_{\min} = 80$, $d = 47$, and
      $80 \bmod 47 = 33 \neq 0$.  But
      $\lfloor 80/47 \rfloor = 1 \neq 2 = \lfloor 106/47 \rfloor$.
      Enumeration reveals $N_0(47) = 1$ (the phantom).
    \item $k = 5$: $\lfloor 242/13 \rfloor = 18 \neq
      19 = \lfloor 250/13 \rfloor$.  Enumeration gives $N_0(13) = 0$.
  \end{itemize}
\end{remark}

% ============================================================
\section{Asymptotic closure via Baker's theorem}\label{sec:baker}
% ============================================================

The Range Exclusion check~\eqref{eq:re} can fail in exactly
two ways.  We show that neither occurs for $k$ sufficiently large.

\subsection{Condition A: floor quotient equality}
\label{ssec:condition-a}

\begin{proposition}\label{prop:cond-a}
  There exists an effectively computable $K_A$ such that for
  all $k > K_A$:
  \[
    \left\lfloor \frac{\corrSum_{\max}}{d} \right\rfloor
    = \left\lfloor \frac{\corrSum_{\min}}{d} \right\rfloor.
  \]
\end{proposition}

\begin{proof}
  The floor quotients differ only if $\Delta \geq d$, i.e.,
  $3^r - 1 \geq 2^S - 3^k$.  Write $\Lambda = S \ln 2 - k \ln 3$.
  By the theorem of Laurent--Mignotte--Nesterenko~\cite{LMN1995}
  on linear forms in two logarithms:
  \begin{equation}\label{eq:lmn}
    |\Lambda| > \exp(-C),
  \end{equation}
  where $C$ is an effectively computable constant depending only on
  the logarithmic heights of~$2$ and~$3$.
  Using the conservative normalization
  ($h_1 = 1$, $h_2 = \ln 3$), we obtain
  $C \leq 24.34 \cdot 1 \cdot \ln 3 \cdot 21^2 \approx 11793$.

  Since $d = 3^k(e^{|\Lambda|} - 1) \geq 3^k \cdot e^{-C}$ and
  $\Delta < 3^r \leq 3^{0.585k+1}$, the condition $\Delta < d$ holds
  whenever
  $3^{0.585k+1} < 3^k \cdot e^{-C}$, i.e.,
  $0.415k - 1 > C/\ln 3$, giving
  \[
    K_A = \left\lceil \frac{C}{0.415 \cdot \ln 3} + \frac{1}{0.415}
    \right\rceil \approx 25868.
  \]
  We take $K_A = 26000$ conservatively.
\end{proof}

\subsection{Condition B: non-divisibility}\label{ssec:condition-b}

The second condition in~\eqref{eq:re} can fail if $d \mid (3^k - 1)$,
i.e., $d \mid \corrSum_{\min}$.

\begin{proposition}\label{prop:cond-b}
  For each fixed positive integer~$q$, the equation
  \begin{equation}\label{eq:pillai}
    (q+1) \cdot 3^k = q \cdot 2^S + 1,
    \quad S = \lceil k \log_2 3 \rceil,
  \end{equation}
  has at most finitely many solutions in~$k$, with an effectively
  computable upper bound on~$k$ depending on~$q$.
\end{proposition}

\begin{proof}
  Equation~\eqref{eq:pillai} is a Pillai-type exponential
  Diophantine equation in the two unknowns $k$ and~$S$.
  Since $S = \lceil k \log_2 3 \rceil$, the pair $(k, S)$ is
  determined by~$k$.  Rewriting:
  $\ln((q+1) \cdot 3^k) - \ln(q \cdot 2^S + 1) = 0$, or equivalently,
  $|S \ln 2 - k \ln 3 - \ln((q+1)/q)| < c/2^S$ for some constant~$c$.
  By Baker's theorem~\cite{BakerWustholz1993}, the left-hand side
  is either~$0$ or bounded below by $\exp(-C'(\log k)^2)$ for an
  effective~$C'$.  Since $c/2^S$ decays exponentially in~$k$, there
  are at most finitely many solutions.
\end{proof}

\begin{remark}\label{rem:cond-b-verification}
  The condition $d \nmid (3^k - 1)$ has been verified
  by exact integer arithmetic for all $k \in \{6, \ldots, 50000\}$.
  No violation has been found.  The Pillai-type argument of
  \cref{prop:cond-b} provides the theoretical guarantee for
  $k \to \infty$, but the effective bound on~$q$ is not
  explicitly computed in this paper.
\end{remark}

\subsection{Combining conditions A and B}\label{ssec:combining}

\begin{proof}[Proof of \cref{thm:main}]
  We distinguish five regimes:
  \begin{enumerate}[label=\textbf{(\roman*)}]
    \item $k = 3, 5$: Direct enumeration shows $N_0(d(k)) = 0$
          (Lean verified).
    \item $k = 4$: $N_0(d(4)) = 1$ (phantom). No actual cycle
          by Simons--de~Weger~\cite{SimonsDeWeger2005}.
    \item $k = 6, \ldots, 10000$: Range Exclusion verified by
          Lean~4 (\cref{thm:finite}).
    \item $k = 10001, \ldots, 50000$: Range Exclusion verified
          by Python exact arithmetic.
    \item $k > 50000$: By \cref{prop:cond-a}, floor quotients
          are equal for $k > K_A = 26000$. By \cref{prop:cond-b}
          and the verification up to $k = 50000$, condition~B
          holds. Hence Range Exclusion holds. \qedhere
  \end{enumerate}
\end{proof}

% ============================================================
\section{Lean~4 formalization}\label{sec:lean}
% ============================================================

The core argument (regimes (i)--(iii) of the proof of
\cref{thm:main}) is formalized in Lean~4 (version~4.28.0).
The formalization resides in the file
\texttt{RangeExclusion.lean} and contains:

\begin{itemize}
  \item Definitions of $S(k)$, $d(k)$, $\corrSum_{\min}(k)$,
    $\corrSum_{\max}(k)$, and the \texttt{checkRE} function
    implementing~\eqref{eq:re}.
  \item A \texttt{checkOne} function that dispatches to
    enumeration for $k \in \{3, 5\}$, skips $k = 4$ (phantom),
    and applies Range Exclusion for $k \geq 6$.
  \item Ten batch verification theorems, each proved by
    \texttt{native\_decide}:
    \begin{center}
    \begin{tabular}{ll}
      \texttt{verify\_3\_to\_1000} & $k = 3, \ldots, 1000$ \\
      \texttt{verify\_1001\_to\_2000} & $k = 1001, \ldots, 2000$ \\
      $\vdots$ & $\vdots$ \\
      \texttt{verify\_9001\_to\_10000} & $k = 9001, \ldots, 10000$ \\
    \end{tabular}
    \end{center}
  \item A combined certificate:
    \texttt{no\_nontrivial\_cycle\_certificate :
    checkRange 3 10000 = true}.
\end{itemize}

The formalization has:
\begin{itemize}
  \item Zero \texttt{sorry} (no unproved obligations).
  \item Two \textbf{axioms}:
    \begin{enumerate}
      \item \texttt{simons\_de\_weger}: no positive cycle with
        $k < 68$ (referencing~\cite{SimonsDeWeger2005}).
      \item \texttt{baker\_lmn}: for $k \geq 10001$,
        \texttt{checkRE k = true} (referencing~\cite{LMN1995}).
    \end{enumerate}
  \item $9997$ values verified by \texttt{native\_decide}.
  \item Total build time: $337$~seconds on Apple~M1~Pro.
\end{itemize}

\begin{remark}[Nature of \texttt{native\_decide}]\label{rem:native}
  The \texttt{native\_decide} tactic compiles the decision
  procedure to native C code via Lean's compiler, then executes it
  and checks the result against the kernel.  This is \emph{not}
  a heuristic: the kernel verifies the proof term.  The computation
  involves integers with up to~$4771$ digits ($3^{10000}$).
\end{remark}

% ============================================================
\section{Discussion}\label{sec:discussion}
% ============================================================

\subsection{Strengths and novelty}\label{ssec:strengths}

The main contribution is a \emph{uniform} argument
(Range Exclusion) that reduces the cycle exclusion problem
to an arithmetic check on the modulus~$d(k)$.  The argument
is elementary (no character sums, no Fourier analysis, no
algebraic number theory beyond Baker's theorem) and mechanically
verifiable.  The Lean formalization provides a machine-checked
certificate for $9997$ values of~$k$.

\subsection{Limitations and honest assessment}\label{ssec:limits}

We document three limitations:

\begin{enumerate}
  \item \textbf{Baker gap.}
    The Lean verification reaches $k = 10000$, while the
    Baker threshold $K_A \approx 26000$.  The gap
    $k \in [10001, 26000]$ is covered by Python (not Lean).
    Extending the Lean verification to $k = 26000$ would
    eliminate this gap but require significant compilation time.

  \item \textbf{Condition~B (non-divisibility).}
    The statement that $d(k) \nmid (3^k - 1)$ for all~$k$
    relies on \cref{prop:cond-b} (Pillai-type finiteness)
    combined with verification up to $k = 50000$.  The effective
    bound from Baker's theorem is not explicitly computed; a
    determined skeptic could demand this computation.  We note
    that no violation has been found computationally, and the
    Pillai-type argument is standard in the theory of exponential
    Diophantine equations.

  \item \textbf{Upper bound (Schur-convexity).}
    The proof that $\corrSum_{\max} = 3^k + 3^r - 2$ uses
    Schur-convexity of the weighted exponential sum.  This has
    been verified by enumeration for $k \leq 20$ but not formally
    proved in Lean.  A formal proof would use the convexity of
    $x \mapsto 2^x$ and the theory of majorization.
\end{enumerate}

\begin{remark}[On the scope of the result]\label{rem:scope-discussion}
  \Cref{thm:main} excludes nontrivial \emph{positive} cycles.
  It does not address the full Collatz conjecture, which also
  requires showing that all orbits reach~$1$ (i.e., excluding
  divergent trajectories).  The almost-all result of
  Tao~\cite{Tao2022} addresses the latter but through entirely
  different methods.
\end{remark}

\subsection{Reproducibility}\label{ssec:reproducibility}

All code (Lean~4 formalization, Python verification script) is
publicly available at:
\begin{center}
  \url{https://github.com/ericmerle3789/Collatz-Junction-Theorem}
\end{center}
Branch: \texttt{proof-assembly-v1}.
The Python script \texttt{verify\_all\_k.py} runs in approximately
$52$~seconds on a modern machine and verifies Range Exclusion for
$k = 6, \ldots, 50000$ using exact integer arithmetic (no floating
point).

\subsection{Open questions}\label{ssec:open}

\begin{enumerate}
  \item Can the Baker constant~$C$ be computed explicitly enough
        to bring $K_A$ down to $K_A \leq 10000$, matching the
        Lean verification range?
  \item Can Condition~B ($d \nmid (3^k-1)$) be proved unconditionally
        for all~$k$, without relying on the Pillai-type finiteness
        argument?
  \item Does Range Exclusion extend to \emph{negative} cycles or
        to generalized $qx{+}r$ dynamics?
  \item Can the Schur-convexity of the upper bound be formalized
        in Lean~4 using Mathlib?
\end{enumerate}

% ============================================================
% Acknowledgements
% ============================================================
\section*{Acknowledgements}

The author thanks the developers of Lean~4 and Mathlib for
providing the infrastructure for formal verification.
Computational work was performed on an Apple M1 Pro (16\,GB).
This research received no external funding.

% ============================================================
% References
% ============================================================
\begin{thebibliography}{99}

\bibitem{BakerWustholz1993}
A.~Baker and G.~W\"ustholz,
\emph{Logarithmic forms and group varieties},
J. reine angew. Math. \textbf{442} (1993), 19--62.

\bibitem{Barina2021}
D.~Barina,
\emph{Convergence verification of the Collatz problem},
J. Supercomput. \textbf{77} (2021), 2681--2688.

\bibitem{BoehmSontacchi1978}
C.~B\"ohm and G.~Sontacchi,
\emph{On the existence of cycles of given length in integer
sequences like $x_{n+1} = x_n/2$ if $x_n$ even, and
$x_{n+1} = 3x_n + 1$ otherwise},
Atti Accad. Naz. Lincei Rend. Cl. Sci. Fis. Mat. Nat.
\textbf{64} (1978), 260--264.

\bibitem{Crandall1978}
R.~E.~Crandall,
\emph{On the ``$3x + 1$'' problem},
Math. Comp. \textbf{32} (1978), 1281--1292.

\bibitem{Eliahou1993}
S.~Eliahou,
\emph{The $3x + 1$ problem: new lower bounds on nontrivial
cycle lengths},
Discrete Math. \textbf{118} (1993), 45--56.

\bibitem{Gouillon2006}
N.~Gouillon,
\emph{Explicit lower bounds for linear forms in two logarithms},
J. Th\'eorie Nombres Bordeaux \textbf{18} (2006), 125--146.

\bibitem{Korec1994}
I.~Korec,
\emph{A density estimate for the $3x + 1$ problem},
Math. Slovaca \textbf{44} (1994), 85--89.

\bibitem{Lagarias1985}
J.~C.~Lagarias,
\emph{The $3x + 1$ problem and its generalizations},
Amer. Math. Monthly \textbf{92} (1985), 3--23.

\bibitem{Lagarias2010}
J.~C.~Lagarias (ed.),
\emph{The Ultimate Challenge: The $3x + 1$ Problem},
AMS, 2010.

\bibitem{LMN1995}
M.~Laurent, M.~Mignotte, and Y.~Nesterenko,
\emph{Formes lin\'eaires en deux logarithmes et d\'eterminants
d'interpolation},
J. Number Theory \textbf{55} (1995), 285--321.

\bibitem{Salikhov2007}
V.~Kh.~Salikhov,
\emph{On the irrationality measure of $\ln 3$},
Dokl. Akad. Nauk \textbf{417} (2007), 753--755.

\bibitem{SimonsDeWeger2005}
D.~Simons and B.~de~Weger,
\emph{Theoretical and computational bounds for $m$-cycles of
the $3n + 1$ problem},
Acta Arith. \textbf{117} (2005), 51--70.

\bibitem{Steiner1977}
R.~P.~Steiner,
\emph{A theorem on the Syracuse problem},
Proc. 7th Manitoba Conf. Numer. Math. Comput., 1977, 553--559.

\bibitem{Tao2022}
T.~Tao,
\emph{Almost all orbits of the Collatz map attain almost bounded
values},
Forum Math. Pi \textbf{10} (2022), e12.

\bibitem{Wirsching1998}
G.~J.~Wirsching,
\emph{The Dynamical System Generated by the $3n + 1$ Function},
Lecture Notes in Mathematics \textbf{1681}, Springer, 1998.

\bibitem{WuWang2014}
Q.~Wu and L.~Wang,
\emph{On the irrationality measure of $\log 3$},
J. Number Theory \textbf{142} (2014), 264--273.

\end{thebibliography}

\end{document}
